%% ma: L12-B14-0078
%% lop: 12
%% chuong: 5
%% bai: 14
%% dang: 12.14.22
%% loai: TN4
%% muc: TH
%% dap_an: "A"
%% nguon: "To_n_12_PT_M_t_Ph_ng.pdf (Đề kiểm tra đánh giá kiến thức, Chuyên đề Phương trình mặt phẳng, Nguyễn Anh Khôi - Lê Minh Khôi), Câu 33"
%% kien_thuc: "Bài 14 (điều kiện hai mặt phẳng vuông góc, điểm thuộc mặt phẳng)"
%% trang_thai: cho-duyet
%% ly_do: "Dạng 12.14.22 mới đề xuất, chờ giáo viên duyệt."
%% ghi_chu: "Tài liệu gốc không kèm đáp án; đáp án do Claude giải và kiểm bằng sympy. Ý A cần giải hệ ba phương trình bậc nhất ba ẩn (máy tính cầm tay được phép)."
\begin{ex}
Trong không gian $Oxyz$, cho ba mặt phẳng $(P)\colon 2x+y+z+3=0$, $(Q)\colon x-y-z-1=0$, $(R)\colon y-z+2=0$. Trong các mệnh đề sau, mệnh đề nào sai?
\choice{\True Không có điểm nào cùng thuộc ba mặt phẳng trên.}{$(P)\perp(R)$}{$(Q)\perp(R)$}{$(P)\perp(Q)$}
\loigiai{$\vec n_P=(2;1;1)$, $\vec n_Q=(1;-1;-1)$, $\vec n_R=(0;1;-1)$. $\vec n_P\cdot\vec n_R=0$, $\vec n_Q\cdot\vec n_R=0$, $\vec n_P\cdot\vec n_Q=0$ nên B, C, D đúng. Giải hệ ba phương trình được nghiệm duy nhất $\left(-\dfrac23;-\dfrac{11}6;\dfrac16\right)$, nên ba mặt phẳng có điểm chung, A sai. Chọn A.}
\end{ex}
